\section{El modelo de competencia espacial: desarrollo y fuentes}
\label{anexo:modelo}

Este anexo desarrolla el modelo que la \Cref{prop:margen} y el
\Cref{cor:entrada} resumen en el cuerpo. Ninguno de los resultados es original:
todos provienen de \textcite{SALOP1979} y el anexo indica, para cada uno, la
ecuación y la página del artículo de las que se obtiene.\footnote{El artículo se
publicó en el \emph{Bell Journal of Economics}, vol.~10, n.º~1, pp.~141--156.
Conviene advertir que Salop no enuncia sus resultados como proposiciones ni
teoremas numerados, sino como ecuaciones dentro del desarrollo del texto; la
numeración empleada aquí y en el cuerpo es propia de este trabajo y sirve
únicamente para referirse a ellos.} El propósito del anexo es doble: dejar
trazable el origen de cada expresión y reunir las limitaciones del marco, que en
el cuerpo se omiten por brevedad.

\subsection{El problema del consumidor}

El mercado local se representa por una circunferencia de perímetro unitario
sobre la que se distribuyen uniformemente $L$ consumidores, cada uno con demanda
unitaria e inelástica por combustible. Operan $n$ estaciones equiespaciadas, de
modo que la distancia entre dos contiguas es $1/n$. Un consumidor localizado a
distancia $x$ de la estación $i$ obtiene un excedente
\begin{equation}
    U_i(x) = v - p_i - c\,x ,
    \label{eq:anexo-utilidad}
\end{equation}
donde $v$ es el precio de reserva efectivo —la valoración del combustible neta
del excedente que otorga el bien externo, ecuación (4) de
\textcite[142]{SALOP1979}— y $c>0$ el costo de desplazamiento por unidad de
distancia. El consumidor compra una unidad en la estación que le entregue el
mayor excedente, siempre que este sea no negativo.

La formulación \eqref{eq:anexo-utilidad} recoge dos rasgos del mercado chileno
descritos en el cuerpo. Primero, la homogeneidad del producto: dos EDS que cobran
el mismo precio son indiferentes para el consumidor, y solo la distancia las
separa. Segundo, la observabilidad de precios, que permite tratar al consumidor
como informado sobre $p_i$ en todo el mercado local.

\subsection{Las tres regiones de la curva de demanda}

\textcite[143--145]{SALOP1979} distingue tres regiones según cuán agresivo sea el
precio de la estación representativa en relación con sus vecinas, ilustradas en
su Figura~1.

\paragraph{Región de monopolio local.} Si el precio de $i$ es suficientemente
alto, la estación solo atrae a los consumidores para quienes el excedente
\eqref{eq:anexo-utilidad} es no negativo, sin llegar a disputar clientes a sus
vecinas. El consumidor marginal se ubica a una distancia $\hat{x}=(v-p_i)/c$, de
modo que la demanda es
\begin{equation}
    q^{m} = \frac{2L}{c}\,(v - p_i) ,
    \label{eq:anexo-qm}
\end{equation}
que es la ecuación (6) de \textcite[143]{SALOP1979}. En esta región la estación
compite contra el bien externo —no comprar, o abastecerse fuera del mercado
local— y no contra sus vecinas.

\paragraph{Región competitiva.} Si el precio es lo bastante bajo como para que
los mercados potenciales de estaciones contiguas se traslapen, el consumidor
marginal ya no es aquel indiferente entre comprar y no comprar, sino aquel
indiferente entre dos estaciones. La demanda es entonces
\begin{equation}
    q^{c} = \frac{L}{c}\left(\bar{p} + \frac{c}{n} - p_i\right) ,
    \label{eq:anexo-qc}
\end{equation}
donde $\bar{p}$ es el precio de las vecinas; es la ecuación (9) de
\textcite[144]{SALOP1979}. Esta es la región relevante para el análisis, pues es
la que corresponde a los mercados urbanos donde ocurre la mayoría de las
entradas.

\paragraph{Región supracompetitiva.} Con precios aún menores la estación captura
la totalidad del mercado de su vecina, situación que aparece bajo la ecuación
(12) de \textcite[145]{SALOP1979}. No se considera aquí, pues corresponde a
comportamientos predatorios ajenos al fenómeno estudiado.

Derivando \eqref{eq:anexo-qm} y \eqref{eq:anexo-qc} se obtienen las pendientes de
la demanda en cada región, $-2L/c$ y $-L/c$ respectivamente, que son las
ecuaciones (10) y (11) de \textcite[144]{SALOP1979}. La distinción entre las dos
primeras regiones tiene una consecuencia empírica directa, que la
\Cref{hip:distancia} recoge y que la \Cref{prop:monopolio} formaliza más abajo.

\subsection{El margen de equilibrio}

Cada estación elige $p_i$ tomando como dados los precios de sus vecinas y
maximizando $(p_i - m)\,q_i - F$, donde $m$ es el costo marginal, común a todas
las estaciones por el argumento del MEPCO expuesto en el cuerpo, y $F$ el costo
fijo de operación del sitio. Imponiendo la condición de primer orden —ecuación
(13) de \textcite[147]{SALOP1979}— sobre la pendiente de la demanda competitiva
se obtiene el resultado de la \Cref{prop:margen},
\begin{equation}
    p = m + \frac{c}{n} ,
    \qquad\text{esto es}\qquad
    \mu = \frac{c}{n} ,
    \label{eq:anexo-margen}
\end{equation}
que es la ecuación (18) de \textcite[147]{SALOP1979}. El margen es proporcional
al costo de desplazamiento e inversamente proporcional al número de competidoras;
equivalentemente, y puesto que las estaciones están equiespaciadas, es
proporcional a la distancia $1/n$ que separa a la estación de sus vecinas. Esta
doble lectura es la que permite que el margen opere en la estrategia empírica
como medida de aislamiento espacial.

\subsection{El umbral: monopolio local}

\begin{proposicion}[Monopolio local]
\label{prop:monopolio}
Si los mercados de estaciones contiguas no se traslapan, la estación resuelve un
problema de monopolio frente al bien externo y su precio óptimo es
\begin{equation}
    p^{m} = \frac{v + m}{2} ,
    \label{eq:monopolio}
\end{equation}
que no depende de $n$ ni de la ubicación de las vecinas.
\end{proposicion}

Se obtiene maximizando el beneficio de monopolio de la ecuación (28) de
\textcite[151]{SALOP1979}, cuya solución es la cantidad $q^{m}=(L/c)(v-m)$ de su
ecuación (29), y sustituyéndola en la demanda de monopolio
\eqref{eq:anexo-qm}, que invertida es $p = v - (c/2L)\,q$. Es el único de los
resultados empleados que Salop no escribe en términos de precio, pues su interés
en esa sección está en las condiciones de rentabilidad y no en el precio de
monopolio.

La proposición es la que dota al modelo de un umbral. Mientras no haya traslape
de mercados, la estructura del mercado local es irrelevante para el precio, de
modo que una entrada lo bastante lejana deja a la incumbente exactamente donde
estaba. El modelo no predice, a nivel de estación, una atenuación asintótica sino
la existencia de un umbral.

\subsection{Entrada libre}

\begin{proposicion}[Entrada libre]
\label{prop:libre}
Bajo libre entrada, el número de estaciones y el margen de equilibrio son
\begin{equation}
    n^{*} = \sqrt{cL/F} ,
    \qquad
    \mu^{*} = \frac{c}{n^{*}} = \sqrt{cF/L} ,
    \label{eq:entrada-libre}
\end{equation}
de modo que el número de estaciones es creciente en la densidad de consumidores
del mercado local y decreciente en el costo fijo de instalación, y el margen se
mueve en sentido inverso.
\end{proposicion}

Son las ecuaciones (19), p.~147, y (22), p.~148, de \textcite{SALOP1979}, que el
autor obtiene de la condición de beneficio nulo (14) evaluada en el equilibrio
competitivo; la expresión para $n^{*}$ se repite como su ecuación (23). Mercados
más densos sostienen más estaciones y márgenes menores.

Conviene señalar que estas dos predicciones —la \Cref{prop:margen} y la
\Cref{prop:libre}— ya han sido contrastadas de manera directa en este mismo
mercado. \textcite{CLEMENZGUGLER2006} toman el modelo de \textcite{SALOP1979}
como marco teórico explícito y derivan de él dos hipótesis: que las estaciones se
ubican más densamente donde la densidad de población es mayor, y que el precio de
equilibrio es menor mientras mayor sea la densidad de estaciones. Encuentran
apoyo para ambas en el mercado minorista austríaco. La densidad de población
explica más del $95\%$ de la variación entre distritos en la densidad de
estaciones, y el coeficiente de la densidad de estaciones sobre el precio tiene el
signo predicho y es significativo al $5\%$ o mejor en todas sus especificaciones;
estimar el sistema como ecuaciones simultáneas no altera sus conclusiones y
sugiere que la causalidad corre desde la densidad hacia el precio, que es la
dirección que este trabajo explota usando una entrada como fuente de variación.

Dos detalles de ese trabajo importan aquí. El primero es que predicen que la
densidad de estaciones crece de manera \emph{cóncava} en la densidad de
población, porque la mayor proximidad reduce el precio de equilibrio y con ello
el beneficio que sostiene a cada sitio. Es exactamente la raíz cuadrada de
\eqref{eq:entrada-libre}. El segundo es la lectura de competencia que le dan a su
segunda hipótesis: la ausencia de una relación negativa entre densidad y precio
sería señal de que las estaciones coordinan sus precios en lugar de competir. Que
en este trabajo una entrada reduzca el margen de las incumbentes es, bajo esa
lectura, evidencia de que la competencia local opera y no está suprimida por
coordinación.

Conviene además situar el ejercicio en la tradición que inaugura
\textcite{BRESNAHAN1991}, que infiere la intensidad de la competencia a partir de
la relación entre el número de firmas y el tamaño del mercado en mercados locales
geográficamente aislados. Su hallazgo central es que la conducta competitiva
cambia rápido con las primeras entradas y casi deja de cambiar después: en
mercados con cinco firmas o menos, casi toda la variación ocurre con la entrada
de la segunda o la tercera, y una vez que el mercado tiene entre tres y cinco
firmas la siguiente entrante casi no altera la conducta. Es el patrón que
\eqref{eq:anexo-efecto} predice, puesto que $c/(n(n+1))$ decrece rápidamente en
$n$, y es el mismo mecanismo que sostiene la \Cref{hip:margen}: la incumbente que
más pierde ante una entrada cercana es la que operaba con pocas competidoras y,
por lo tanto, con márgenes altos.

\subsection{El efecto de una entrada y las tres hipótesis}

El \Cref{cor:entrada} es la estática comparativa directa de
\eqref{eq:anexo-margen} respecto de $n$, tratando el número de competidoras como
un parámetro de corto plazo,
\begin{equation}
    \Delta\mu = \frac{c}{n+1} - \frac{c}{n} = -\,\frac{c}{n(n+1)} \;<\; 0 .
    \label{eq:anexo-efecto}
\end{equation}
No aparece como tal en \textcite{SALOP1979}, cuya sección 4 (pp.~148--149)
desarrolla las estáticas comparativas respecto de los parámetros
$\{F, m, v, c, L\}$ y no respecto del número de firmas, que en su análisis de
largo plazo es endógeno. La diferencia es de horizonte y no de contenido: la
\Cref{prop:libre} describe hacia dónde converge el mercado en el largo plazo,
mientras que \eqref{eq:anexo-efecto} describe el ajuste de precios de las
incumbentes ante un cambio en $n$ que ellas toman como dado, que es el objeto de
este trabajo.

Dado que $m$ está determinado por una referencia nacional y la entrada es un
evento local, el costo marginal no se altera y la caída del margen se traslada
uno a uno al precio, $\Delta p = \Delta\mu$. Esta equivalencia es la que permite
contrastar el modelo tanto con precios como con márgenes, y es consecuencia
directa del carácter nacional del MEPCO.

\paragraph{Sobre el nivel (\Cref{hip:nivel}).} El signo de
\eqref{eq:anexo-efecto} es inequívoco para cualquier entrada que aumente el
número de competidoras efectivas de la incumbente. El efecto es además
permanente, pues opera sobre el equilibrio del mercado local y no sobre una
desviación transitoria respecto de él.

\paragraph{Sobre la distancia (\Cref{hip:distancia}).} Por la
\Cref{prop:monopolio}, una entrada que no genere traslape de mercados deja a la
incumbente frente al bien externo, donde su precio óptimo \eqref{eq:monopolio} es
independiente de la estructura del mercado local, y el efecto es exactamente
nulo. Conviene precisar qué predice el modelo y qué no. Para una estación dada la
predicción es discontinua: la entrante cuenta en $n_i$ o no cuenta, y el efecto es
$-c/(n_i(n_i+1))$ o cero. Lo que se estima empíricamente, en cambio, es el efecto
promedio por anillo de distancia sobre estaciones cuyo radio de mercado difiere.
Como la proporción de incumbentes para las cuales la entrante cae dentro del
radio decrece con la distancia, el promedio de umbrales heterogéneos genera un
perfil que decae de manera suave y se anula más allá de cierta distancia. Es ese
perfil, y no una atenuación asintótica a nivel de estación, lo que la
\Cref{hip:distancia} afirma y lo que la \Cref{fig:atenuacion} contrasta.

\paragraph{Sobre el margen previo (\Cref{hip:margen}).} Por
\eqref{eq:anexo-margen}, el margen es un estadístico suficiente del aislamiento
espacial de la estación: $\mu$ es alto precisamente cuando las competidoras son
pocas y están lejos. Invirtiendo \eqref{eq:anexo-margen} como $n = c/\mu$ y
sustituyendo en \eqref{eq:anexo-efecto}, el efecto de una entrada se expresa
directamente en función del margen previo,
\begin{equation}
    |\Delta\mu| = \frac{\mu^{2}}{c + \mu} ,
    \qquad
    \frac{\partial\,|\Delta\mu|}{\partial\mu}
    = \frac{\mu\,(2c + \mu)}{(c + \mu)^{2}} \;>\; 0 ,
    \label{eq:het}
\end{equation}
de modo que una incumbente con margen previo elevado es una incumbente aislada y
es, por lo tanto, la que más pierde ante una entrada cercana. La predicción vale
además en términos relativos, que es la escala en que se estiman los resultados
sobre el precio en logaritmo: $|\Delta\mu|/\mu = \mu/(c+\mu)$ es también
creciente en $\mu$.

Vale la pena notar que la formulación en términos de $\mu$ y no del número de
competidoras dentro de un radio fijo no es una elección de conveniencia. El
conteo $n_i(r) = \#\{j : \mathrm{dist}(i,j) \le r\}$ se construye sobre un radio
$r$ igual para todas las estaciones, mientras que el $n_i$ del modelo es el
número de competidoras dentro del radio propio de cada mercado local, que no se
observa. Por \eqref{eq:anexo-margen}, en cambio, $\mu$ es una transformación
exacta de ese $n_i$. El margen es, en este modelo, el estadístico suficiente del
poder de mercado local, y el conteo a radio fijo apenas una aproximación gruesa.

\subsection{Limitaciones del marco}

Tres limitaciones conviene declarar. La primera es que el equilibrio de
\textcite{SALOP1979} predice un margen común a todas las estaciones que enfrentan
el mismo $n$, mientras que los datos muestran dispersión de precios persistente
dentro de mercados muy estrechos. Leer la circunferencia como el mercado local de
cada estación, y no como una unidad administrativa, acomoda parte de esa
dispersión, pues estaciones de una misma comuna enfrentan distinto $n_i$; pero no
toda. La literatura atribuye el resto a fricciones informativas del lado del
consumidor \parencite{FISCHER2025}, un canal que este modelo no incorpora.

La segunda es que el equiespaciamiento es un supuesto de conveniencia y no un
rasgo del mercado chileno. Su función es hacer que $1/n$ sea a la vez el
espaciamiento y el inverso del número de competidoras, de modo que ambas lecturas
de \eqref{eq:anexo-margen} coincidan. Con espaciamiento desigual las magnitudes
cambian, pero no el signo ni el ordenamiento en que descansan las tres hipótesis.

La tercera es que la \Cref{prop:libre} describe la entrada como una respuesta a
la densidad del mercado local, es decir, como una decisión endógena respecto de
dónde instalarse. La estrategia de identificación desarrollada en el cuerpo no
requiere negar esa endogeneidad sino únicamente que el momento preciso en que la
entrada se materializa sea exógeno desde la perspectiva de la incumbente.
